\chapter{Conclusions and Future Work}
\label{ch:conclusions}

\section{Conclusions}

This thesis addressed the problem of evaluating pit decision support systems from historical Formula 1 data under live-like temporal constraints, the implemented system
starts from historical race sessions, converts them into replayable event streams, processes them through a Flink based strategy pipeline, and then compares three decision
paradigms under shared evaluation contracts: the Flink Strategy Engine, Batch XGBoost, and the MOA Online Learner.

A pit recommendation is emitted at a specific driver-lap, while its correctness can only be assessed later,
after the race has evolved and the pit event, if any, has occurred. The evaluation therefore had to connect decision rows, future pit events, and post pit outcome labels
without using future information as model input. This led to the two contracts used throughout the thesis: \texttt{pit\_any\_h2}, for pit timing, and \texttt{pit\_success\_h2},
for successful pit-call prediction.

The final evaluation showed that these two contracts describe different levels of difficulty, where Pit timing can be learned from replayed race context with useful precision
and coverage trade-offs, while Successful pit-call prediction can also be evaluated under a strict operational view, but it remains sparse, threshold sensitive, and harder
to cover reliably. The following findings summarize the main outcomes of the work.

\subsection{Principal Findings}

\begin{enumerate}
    \item \textbf{Aligned contracts are necessary for a fair comparison.}

          The experiments showed that the three systems cannot be compared directly from their native outputs. The Flink Strategy Engine emits rule based labels, Batch
          XGBoost emits out-of-fold scores, and MOA emits vote based scores in prequential order. A meaningful comparison required a shared truth universe, explicit truth
          lenses, leakage filtering, and a clear separation between row-level decisions and event-level coverage. These elements are part of the contribution of the thesis,
          because they define the conditions under which deterministic rules, Batch learning, and online learning can be judged on the same evidence.

    \item \textbf{\texttt{pit\_any\_h2} is the cleaner learnable task.}

          The pit timing contract gave the clearest learning signal. It checks whether an eligible pit stop happens within the evaluation horizon, without asking whether
          that pit is later classified as strategically successful. Under this contract, the learning based systems improved over the deterministic Flink Strategy Engine.
          Batch Percent reached the largest event recall, while MOA gave the strongest precision oriented balance. The task remains imbalanced, but the results indicate
          that race context features such as pace, tyre life, gaps, and race progress contain usable signal for predicting that a pit window is approaching.

    \item \textbf{\texttt{pit\_success\_h2} is stricter and policy-sensitive.}

          The success contract was substantially harder. A positive call is correct only when it matches a future eligible pit and the post pit evaluator classifies that
          pit as successful. This reduces the amount of clean positive evidence and makes the selected operating threshold much more influential. MOA Percent produced the
          cleanest selected successful pit calls, but its successful event coverage remained low. The relaxed threshold analysis confirmed the expected trade-off: coverage
          can be increased, but only by accepting more false positives under the strict operational view.

    \item \textbf{The Flink Strategy Engine remains useful as an interpretable baseline.}

          The rule based system did not dominate the strict precision results, but it provided an interpretable reference point. Its decisions can be followed through explicit
          score terms, gates, and promotion rules. The evaluation also showed why diagnostic matched pit quality must be separated from operational precision. When a Flink call
          matches a known pit, many of those matched pits correspond to successful outcomes. Once no-match calls and failed matched pits are counted as false positives, strict
          precision becomes much lower. This distinction prevents the deterministic baseline from being overstated and keeps the comparison with Batch and MOA fair.

\end{enumerate}

Taken together, these findings answer the two main contract questions of the thesis; pit timing is a learnable decision support task in the implemented replay setting,
successful pit-call prediction can be evaluated under a stricter contract, but the results show a more limited and fragile signal. The strongest conclusion is 
that each paradigm occupies a different position in the trade-off between interpretability, coverage, precision, and temporal evaluation.

\section{Scope and Limitations}

The results should be read within the scope of the implemented pipeline; the first boundary is the use of historical replay, replaying FastF1 sessions through Kafka and
Flink makes the experiments reproducible and gives the system live-like temporal semantics, the same race can be replayed multiple times, the produced artifacts can be
inspected, and the comparison can be regenerated. A live race feed would introduce additional issues, such as data corrections, missing fields, feed latency, and operational
failure modes. The current system is therefore a controlled experimental environment rather than a deployed pit-wall tool usable in production for real world applications.

A second boundary concerns the definition of successful pit outcomes, the \texttt{pit\_success\_h2} target depends on the \texttt{PitStopEvaluator}, which classifies completed
pit cycles from post pit evidence. This gives the thesis a reproducible truth contract, but it does not represent the full utility function of a racing team because a real strategy
decision may also depend on tyre allocation, risk tolerance, team objectives, opponent response, long-term stint plans, or championship context. The evaluator captures part of
the strategic outcome, mainly through gap and rival evolution, but it remains an heuristic approximation.

The strict success contract also creates a support problem, successful pit events are rarer than pit timing events, and some races contain only a small number of clean
positive cases. Race-level Kappa and G-Mean diagnostics therefore become noisy for \texttt{pit\_success\_h2}, even after support filtering. This sparsity limits how strongly
per race behaviour can be interpreted.

Threshold choice is another limitation since Batch XGBoost and MOA produce scores, and the final operational behaviour depends on where the threshold is placed. The selected
thresholds define the headline policies, while the relaxed rows show how quickly the balance between precision and coverage can change. A fixed threshold is useful for
reproducibility, but it cannot represent every possible operational preference.

The MOA branch adds two further boundaries; its positive scores are derived from class votes and are treated as ranking scores, they support
threshold sweeps and precision--recall diagnostics, but they should not be read as direct probabilities of a pit event. The same caution applies to explainability, Batch
XGBoost can be inspected directly through SHAP and TreeSHAP, while the MOA Online Learner does not expose an equally direct attribution path. A surrogate SHAP analysis was
explored, but it was not retained because the surrogate model could not be verified as a faithful explanation of the internal online learner.

Finally, the model comparison is deliberately scoped; the Batch branch uses XGBoost and the online branch uses the selected MOA configuration. This is enough to evaluate the
proposed framework across deterministic, Batch, and online paradigms, but it does not exhaust the space of possible learners. Other Batch models, online learners, calibration
methods, and drift-aware algorithms may produce different trade-offs under the same contracts.

These limitations define what the comparison actually supports: a reproducible evaluation of three decision paradigms on historical
replay data, under explicit contracts and controlled accounting rules. They also indicate where future work should extend the system.

\section{Future Work}

The implemented pipeline provides a foundation for several extensions, the most relevant ones follow directly from the limitations described above.

\subsection{From Historical Replay to Live Ingestion}

The first extension is the transition from replay to live ingestion, the current architecture already separates event sourcing, Kafka streaming, Flink processing,
persisted artifacts, and downstream evaluation. A live version would replace the historical replay source with a real timing or telemetry feed, while keeping the same
stateful processing logic where possible.

This step would test a different part of the system, in replay the input sequence is frozen and can be regenerated while in a live setting, events may arrive late, be
corrected, or be partially unavailable. The system would need stronger monitoring, latency checks, and failure handling, also the evaluation would also have to distinguish
between decisions made with incomplete live data and decisions made after corrected data become available.

\subsection{Richer Truth Models for Successful Pit Outcomes}

The second direction concerns the successful pit truth model. The current \texttt{PitStopEvaluator} provides a deterministic classification of completed pit cycles,
which was necessary for this thesis, a richer evaluator could use longer post pit windows, better rival tracking, tyre compound plans, and more detailed traffic evolution.
It could also include sampled human validation, especially for ambiguous pit cycles where the rule based classification is less clear.

Another possible development is a probabilistic success label, instead of assigning only success, failure, or unknown categories, the evaluator could estimate the
confidence of the outcome; this would help in cases where the gap change is small, the rival context changes after the stop, or safety car and weather conditions
distort the comparison.

\subsection{Event Level and Sequence Aware Learning}

The current learning setup uses driver lap rows, this representation is simple and works well with tabular learners, but part of the evaluation is event level. One future
direction is to train models more directly for event coverage, rather than treating each driver lap row as an independent binary example.

Several formulations are possible, for instance the model could predict time to pit, rank pit opportunities within a short horizon, or optimize event coverage under a false positive
budget. Sequence aware models could also use recent lap history directly, instead of relying only on the engineered state available in the current row. This direction is
especially relevant for \texttt{pit\_success\_h2}, where the signal is sparse and isolated row predictions may not capture the full strategic episode.

\subsection{Adaptive Threshold and Risk Aware Policies}

The threshold frontier showed that the same score model can behave differently depending on the selected operating point, future work could replace fixed thresholds
with adaptive policies. A conservative threshold may be appropriate during stable race phases, while a lower threshold may be useful when a rival has just pitted, when
a safety car appears, or when tyre degradation is rapidly increasing.
Risk-aware policies could expose different operating modes, for instance a monitoring mode could favour coverage and emit early warnings, a recommendation mode could require higher
precision before producing an actionable call. This would make the output more flexible and closer to the way a strategy engineer might use a decision support tool.

\subsection{Broader Online Learning and Calibration Aware Evaluation}

The MOA Online Learner was evaluated through one selected configuration, future work could compare additional online learners, including Hoeffding Trees, Adaptive
Random Forests, and other drift aware ensembles. This would clarify whether the observed MOA behaviour depends on the selected learner or whether the same pattern
appears across different online methods.

\subsection{Faithful Explainability for Online Learners}

The final direction concerns explainability for online models, since the Batch branch has a direct explanation path through tree based attribution, while the MOA branch does not; a
future version could investigate explanation methods designed for streaming learners, or log additional interpretable signals during prequential prediction.
The goal would be to explain online decisions without replacing the learner with a separate surrogate model, possible approaches include local perturbation diagnostics,
model specific explanations when available, or rule summaries extracted from online behaviour. This would make the online branch easier to inspect and more comparable with
the interpretability of the Flink Strategy Engine and Batch XGBoost.

\section{Final Remarks}

The thesis developed a pipeline for comparing pit decision systems from historical race replay to stream processing, learning, and evaluation. The main outcome is a
structured comparison where pit timing showed the clearest learning gains. Successful pit call prediction remained more difficult, mainly
because success is rarer, evaluator dependent, and sensitive to the operating threshold.

The work also showed that evaluation design matters as much as model choice in this setting, a matched pit diagnostic can suggest good conditional behaviour while strict
operational precision remains low. A score can rank rows well while still requiring a conservative threshold. A feature can improve apparent performance while acting as a
shortcut. These effects are part of the pit decision problem and have to be exposed.

The final contribution is therefore both practical and methodological; practically, the system provides an executable path from historical races to comparable decision artifacts,
methodologically, it defines contracts and accounting rules that make deterministic stream rules, Batch learning, and online learning comparable under the same evidence.